\documentclass[11pt, a4paper, spanish]{article}

% ==============================================================================
% PREÁMBULO
% ==============================================================================
\usepackage[utf8]{inputenc}
\usepackage[T1]{fontenc}
\usepackage{babel}
\usepackage{lmodern}
\usepackage{geometry}
\usepackage{graphicx}
\usepackage{enumitem}
\usepackage{hyperref}
\usepackage{booktabs}
\usepackage{tabularx}
\usepackage{xcolor}
\usepackage{titlesec}
\usepackage{array}
\usepackage{ragged2e}

% --- CONFIGURACIONES ---
\geometry{left=2.5cm, right=2.5cm, top=2.5cm, bottom=2.5cm}
\setlength{\parskip}{0.5em}
\hypersetup{colorlinks=true, urlcolor=blue!50!black, linkcolor=blue!50!black}
\renewcommand{\familydefault}{\sfdefault}

% Formato de títulos
\titleformat{\section}{\Large\bfseries\color{blue!50!black}}{\thesection}{1em}{}
\titlespacing*{\section}{0pt}{2ex plus 1ex minus .2ex}{1.5ex plus .2ex}

% Configuración de tablas
\newcolumntype{L}[1]{>{\RaggedRight\arraybackslash}p{#1}}
\newcolumntype{R}[1]{>{\RaggedLeft\arraybackslash}p{#1}}
\newcolumntype{C}[1]{>{\Centering\arraybackslash}p{#1}}

% ==============================================================================
% DOCUMENTO
% ==============================================================================
\begin{document}

% [CONTENIDO PREVIO IGUAL...]

\subsection*{La Hoja de Ruta: De los Datos a la Decisión}
El entregable final fue una hoja de ruta estratégica, diseñada para ser presentada a la dirección, justificando cada acción con datos y un claro Retorno de la Inversión en Seguridad (ROSI).

\begin{table}[h!]
\centering
\footnotesize
\renewcommand{\arraystretch}{1.5}
\begin{tabularx}{\textwidth}{@{}>{\bfseries}L{3cm} L{4.5cm} L{4.5cm} @{}}
    \toprule
    \textbf{Oportunidad Estratégica} & \textbf{Acción de Mitigación Propuesta} & \textbf{Impacto y Justificación (ROSI)} \\
    \midrule
    Mitigación de Riesgos Frecuentes & Implementar \textbf{Zonas de Exclusión Dinámicas} calculadas según la energía potencial de cada izaje (usando un modelo de trayectoria parabólica simple) y reforzar su señalización. & \textbf{Alto Retorno, Baja Inversión.} Mitiga el 85\% de los escenarios de riesgo moderado. Reduce drásticamente la exposición del personal al riesgo diario sin inversión en nuevo equipamiento. \\
    
    \addlinespace
    Control de Riesgos Críticos & Establecer un protocolo de \textbf{Planes de Izaje Críticos (PIC)} para toda maniobra que supere un umbral de energía de 10 kJ, incluyendo revisión obligatoria por supervisión. & \textbf{Control de Procesos Clave.} Sistematiza la gestión de las operaciones más peligrosas, generando un registro auditable y reduciendo el factor de error humano con un costo administrativo marginal. \\
    
    \addlinespace
    Eliminación de Riesgos Catastróficos & Requerir \textbf{sistemas de retención redundantes} (doble eslinga/estrobo) para todo izaje de más de 25 kJ o en áreas con personal no evacuable. & \textbf{Inversión en Resiliencia.} Aunque es un costo mayor, esta medida de ingeniería es la única que mitiga eficazmente el riesgo de falla de equipo, protegiendo vidas y evitando paradas de operación catastróficas. \\
    \bottomrule
\end{tabularx}
\caption{Hoja de ruta estratégica para mitigación de riesgos}
\end{table}

% [RESTO DEL DOCUMENTO IGUAL...]

\end{document}